\section{Метод характеристик. Нехарактеристические граничные условия. Локальная теорема существования.}
\subsection{Метод характеристик}
Пусть область $U \subset \R^n$ имеет гладкую границу $\partial U$, 
а $\gamma \subset \partial U$ — гладкая кривая, на которой зафиксировано 
значение искомой функции. 

Тогда математическая постановка задачи принимает следующий вид:
\begin{equation} \label{eq::1}
\System{
    &F(Du, u, x) = 0, \quad x \in U, \\
    &u\big|_{\gamma} = g,
}
\end{equation}
где $F(p, z, x)$ и $g$ — заданные гладкие функции своих аргументов, 
$p = (p_1, \dots, p_n) \in \R^n$.

Поиск решения можно свести к интегрированию системы обыкновенных дифференциальных 
уравнений вдоль специальных кривых --- характеристик. 
Пусть характеристика задана параметрически как $x(s)$. Обозначим значения 
функции и её градиента на кривой как $z(s) = u(x(s))$ и $p(s) = Du(x(s))$.
    
Зададим уравнение для пространственных координат характеристической кривой $x(s)$ следующим образом:
\begin{equation}
\frac{dx_j}{ds} = \frac{\partial F}{\partial p_j}(p, z, x),
\quad \text{или в векторной форме } \frac{dx}{ds} = D_p F(p, z, x)
\end{equation}

Теперь найдем, как меняется градиент $p_i(s)$ вдоль этой кривой. 
Продифференцируем компоненту $p_i(s) = \frac{\partial u(x(s))}{\partial x_i}$
по параметру $s$, используя правило дифференцирования сложной функции:
\begin{equation} \label{eq::dp_ds_chain}
\frac{dp_i(s)}{ds} = \frac{d}{ds} \left[ \frac{\partial u(x(s))}{\partial x_i} \right] = \sum_{j=1}^n \frac{\partial^2 u(x(s))}{\partial x_j \partial x_i} \frac{dx_j}{ds}
\end{equation}

С другой стороны, поскольку исходное уравнение $F(Du, u, x) = 0$ должно
выполняться тождественно, его полная производная по любой пространственной
переменной $x_i$ равна нулю. Продифференцируем функцию $F(p(x), z(x), x)$:
\begin{equation}
\frac{d}{dx_i} F(p(x), z(x), x) = \sum_{j=1}^n \frac{\partial F}{\partial p_j} \frac{\partial p_j}{\partial x_i} + \frac{\partial F}{\partial z} \frac{\partial z}{\partial x_i} + \frac{\partial F}{\partial x_i} = 0
\end{equation}

Учитывая определения $p_j = \frac{\partial u}{\partial x_j}$, $z = u$, а также $\frac{\partial p_j}{\partial x_i} = \frac{\partial^2 u}{\partial x_j \partial x_i}$ и $p_i = \frac{\partial z}{\partial x_i}$, перепишем это тождество:
\begin{equation} \label{eq::F_diff}
\sum_{j=1}^n \frac{\partial F}{\partial p_j} \frac{\partial^2 u}{\partial x_j \partial x_i} + \frac{\partial F}{\partial z} p_i + \frac{\partial F}{\partial x_i} = 0
\end{equation}

Если мы выберем вектор скорости кривой $\frac{dx_j}{ds} = \frac{\partial F}{\partial p_j}$, то первая сумма в уравнении \eqref{eq::F_diff} в точности совпадет с выражением для $\frac{dp_i(s)}{ds}$ из \eqref{eq::dp_ds_chain}. Подставляя $\frac{dp_i}{ds}$ в \eqref{eq::F_diff}, получаем дифференциальное уравнение для $p_i$:
\begin{equation}
\frac{dp_i(s)}{ds} + \frac{\partial F}{\partial z} p_i + \frac{\partial F}{\partial x_i} = 0 \implies \frac{dp_i(s)}{ds} = -\frac{\partial F}{\partial x_i} - \frac{\partial F}{\partial z} p_i
\end{equation}

Выведем уравнение для изменения значения функции $z(s)$ вдоль характеристики. Дифференцируя $z(s) = u(x(s))$ по $s$, получаем:
\begin{equation}
\frac{dz(s)}{ds} = \frac{d}{ds}(u(x(s))) = \sum_{i=1}^n \frac{\partial u(x(s))}{\partial x_i} \frac{dx_i}{ds}
\end{equation}
Так как $\frac{\partial u}{\partial x_i} = p_i$ и $\frac{dx_i}{ds} = \frac{\partial F}{\partial p_i}$, окончательно имеем:
\begin{equation}
\frac{dz}{ds} = \sum_{i=1}^n p_i \frac{\partial F}{\partial p_i}
\end{equation}

Объединяя все выведенные компоненты, получаем фундаментальную \textit{характеристическую систему} уравнений:
\begin{equation} \label{eq::char_system}
\begin{cases}
\begin{aligned}
\frac{dx_i}{ds} &= F_{p_i}(p, z, x) \\
\frac{dz}{ds} &= \sum_{i=1}^n p_i F_{p_i}(p, z, x) \\
\frac{dp_i}{ds} &= -F_{x_i}(p, z, x) - p_i F_z(p, z, x)
\end{aligned}
\end{cases}
  \quad \text{или} \quad 
\begin{cases}
\begin{aligned}
    \frac{dx}{ds} &= D_p F(p, z, x) \\
    \frac{dz}{ds} &= \langle p, D_p F(p, z, x) \rangle \\
    \frac{dp}{ds} &= -D_x F(p, z, x) - p F_z(p, z, x)
\end{aligned}
\end{cases}
\end{equation}

\begin{theorem}[О структуре характеристических уравнений]
    Если гладкая функция $u(x) \in C^2(U)$ является решением уравнения 
    \eqref{eq::1} $F(Du, u, x) = 0$, 
    а $x(s)$ --- решение системы обыкновенных дифференциальных уравнений вида
    \begin{equation}
        \frac{dx}{ds} = D_p F(p(s), z(s), x(s)),
    \end{equation}
    то тройка функций $(p(s), z(s), x(s))$, где $z(s) = u(x(s))$ и $p(s) = Du(x(s))$, 
    удовлетворяет характеристической системе обыкновенных дифференциальных 
    уравнений \eqref{eq::char_system}.
\end{theorem}

\subsection{Примеры систем характеристик}

\textbf{1. Линейное уравнение}\\
Для уравнения $(b(x), Du) = c(x)u$, функция принимает вид $F(p, z, x) = (b(x), p) - c(x)z$. 
Тогда характеристическая система упрощается до:
\begin{equation}
    \System{
        \frac{dx}{ds} &= b(x) \\
        \frac{dz}{ds} &= c(x)z
    }
\end{equation}

\textbf{2. Квазилинейное уравнение}\\
Рассмотрим уравнение $(b(x, u), Du) + c(x, u) = 0$. 
Здесь $F(p, z, x) = (b(x, z), p) + c(x, z)$. Характеристическая система для пространственных координат и функции принимает вид:
\begin{equation}
    \System{
        \frac{dx}{ds} &= b(x, z) \\
        \frac{dz}{ds} &= -c(x, z)
    }
\end{equation}
Важным свойством квазилинейных уравнений является то, что подсистема для $x$ и $z$ полностью отщепляется от дифференциальных уравнений на $p$, что существенно облегчает её интегрирование.

\subsection{Нехарактеристические граничные условия}

\begin{theorem}[Напоминание: Теорема о неявной функции] \label{th::implicit}
Пусть отображение $G(p, y) \colon \mathbb{R}^n_p \times \mathbb{R}^{n-1}_y \to \mathbb{R}^n$
является непрерывно дифференцируемым ($C^1$) в некоторой окрестности точки
$(p^0, y^0) \in \mathbb{R}^n \times \mathbb{R}^{n-1}$, причем выполнено равенство:
\begin{equation}
  G(p^0, y^0) = 0.
\end{equation}
Если матрица Якоби отображения $G$ по переменным $p$ в точке $(p^0, y^0)$ является 
невырожденной, то есть её якобиан отличен от нуля:
\begin{equation}
\det D_p G(p^0, y^0) \neq 0,
\end{equation}
то существуют открытая окрестность $W \subset \mathbb{R}^{n-1}$ точки $y^0$,
открытая окрестность $V \subset \mathbb{R}^n$ точки $p^0$ и единственное
  отображение $q \colon W \to V$ класса $C^1$, такие что:
\begin{enumerate}
  \item $q(y^0) = p^0$;
  \item для всех $y \in W$ справедливо тождество $G(q(y), y) \equiv 0$.
\end{enumerate}
\end{theorem}

\textbf{1. Выпрямление границы}

Определим область $U = \{y = (y_1, \dots, y_n) \mid y_n \ge \gamma - \varphi(y_1, \dots, y_{n-1})\}$
и ее границу $\Gamma = \{y = (y_1, \dots, y_n) \mid y_n = \gamma - \varphi(y_1, \dots, y_{n-1})\}$.
Здесь предполагается, что функция $\varphi(\cdot)$ принадлежит классу гладкости $C^k$, где $k \ge 2$.
Зафиксируем на границе точку $y^0$ и зададим начальное значение функции 
$z^0 = g(y_1^0, \dots, y_{n-1}^0)$. Компоненты вектора начального градиента
вдоль выпрямленной границы задаются частными производными: $p_i^0 = \frac{\partial \varphi(y_1^0, \dots, y_{n-1}^0)}{\partial y_i}$ для $i = 1, \dots, n-1$.
Недостающую нормальную компоненту $p_n^0$ находим из исходного дифференциального 
уравнения:


$$F(p_1^0, \dots, p_{n-1}^0, p_n^0, z^0, x^0) = 0$$

Согласно теореме о неявной функции \eqref{th::implicit}, для локального разрешения этого уравнения 
относительно $p_n$, производная функции $F$ по этому аргументу должна быть отлична 
от нуля: $\frac{\partial F}{\partial p_n}(p^0, z^0, x^0) \neq 0$. В векторной 
форме это требование эквивалентно тому, что скалярное произведение вектора 
нормали к границе и градиента функции $F$ по переменным $p$ не равно нулю: 
$\langle \nu(x^0), \nabla_p F(p^0, z^0, x^0) \rangle \neq 0$.

\textbf{2. Граничные условия характеристической системы}\\
\begin{equation} \label{eq::2}
\System{
    &F(Du, u, x) = 0, \quad x \in U, \\
    &u\big|_{\Gamma} = g,
}
\end{equation}

Введем диффеоморфизмы $\Phi \colon \mathbb{R}^n \to \mathbb{R}^n$ и обратный к 
нему $\Psi = \Phi^{-1} \colon \mathbb{R}^n \to \mathbb{R}^n$, которые локально 
выпрямляют границу. Рассмотрим функции $u: U\to\R$ и $v(y) = u(\Psi(y))$.
По правилу дифференцирования сложной функции их градиенты связаны как $Du(x) = Dv(y) D\Phi(x)$.
Подставляя это выражение в исходное уравнение, получаем его представление в новых координатах:


$$F(Dv(y)D\Phi(\Psi(y)), v(y), \Psi(y)) = 0$$

\begin{definition}
  Граничное условие называется \textit{нехарактеристическим} в точке 
  $x^0 \in \Gamma$, если вектор $p^0$, найденный из описанной выше 
  процедуры, удовлетворяет условию:
  \begin{equation}
      \left\langle \nu(x^0), \nabla_p F(p^0, z^0, x^0) \right\rangle \neq 0,
  \end{equation}
  где $\nu(x^0)$ --- вектор нормали к поверхности $\Gamma$ в точке $x^0$.
\end{definition}

В случае локально выпрямленной границы $\Gamma = \{x_n = 0\}$, условие 
нехарактеристичности принимает простой вид:
\begin{equation}
    \frac{\partial F}{\partial p_n}(p^0, z^0, x^0) \neq 0.
\end{equation}
Это условие, согласно теореме о неявной функции, гарантирует возможность 
локально разрешить уравнение $F(p^0, z^0, x^0) = 0$ относительно $p_n^0$.




\begin{definition}
Будем говорить, что граница $\partial U$ принадлежит классу $C^k$ ($k \ge 2$),
если для любой точки $x^0 \in \partial U$ существуют число $r > 0$ и функция 
$\gamma \colon \mathbb{R}^{n-1} \to \mathbb{R}$ класса $C^k$ такие, что 
  (при необходимости после поворота и перенумерации координатных осей):
\begin{equation}
    U \cap B(x^0, r) = \left\{ x \in B(x^0, r) \;\middle|\; x_n > \gamma(x_1, \dots, x_{n-1}) \right\},
\end{equation}
% а локальное уравнение самой границы в этом шаре задается как:
% \begin{equation}
%     \partial U \cap B(x^0, r) = \left\{ x \in B(x^0, r) \;\middle|\; x_n = \gamma(x_1, \dots, x_{n-1}) \right\}.
% \end{equation}
\end{definition}

\begin{definition}
Распрямление границы задается следующей заменой координат: 
  $y_i = x_i = \Phi_i(x)$ для $i = 1, \dots, n-1$ и $y_n = x_n - \gamma(x_1, \dots, x_{n-1}) = \Phi_n(x)$.
\end{definition}

\textbf{3. Лемма о характеристических граничных условиях}

\begin{lemma}
Если в точке $x^0$ выполнено условие нехарактеристичности 
$F_{p_n}(p^0, z^0, x^0) \neq 0$ и $F(p^0, z^0, x^0) \neq 0$, то для всех $y \in \Gamma$,
достаточно близких к $x^0$, существует единственное решение системы:
  \begin{equation}
    \begin{cases}
      q_i(y) = \frac{\partial g(y)}{\partial y_i}, \quad i = 1, \dots, n-1, \\
      F(q(y), g(y), y) = 0.
    \end{cases}
  \end{equation}
\end{lemma}
\begin{proof}
Для доказательства разрешимости этой системы относительно неизвестных функций $q(y)$ 
вводится вспомогательное векторное отображение 
$$G(p, y) = (G_1(p, y), \dots, G_n(p, y))^T \colon \mathbb{R}^n \times \mathbb{R}^{n-1} \to \mathbb{R}^n$$
компоненты которого имеют вид:
\begin{equation}
  \begin{cases}
    G_i(p, y) = p_i - \frac{\partial g(y)}{\partial y_i} \quad (i = 1, \dots, n-1),\\
    \quad G_n(p, y) = F(p, g(y), y).
  \end{cases}
\end{equation}
Тогда задача сводится к нахождению корней неявной системы уравнений $G(p^0, x^0) = 0$.

Найдем матрицу Якоби $D_p G(p^0, x^0)$. В первых $n-1$ строках на главной диагонали 
стоят единицы, а последняя строка состоит из частных производных исходного
уравнения $F_{p_1}(p^0, x^0), \dots, F_{p_n}(p^0, x^0)$.
\begin{equation}
D_p G(p^0, x^0) = \begin{pmatrix}
1 & 0 & \dots & 0 & 0 \\
0 & 1 & \dots & 0 & 0 \\
\vdots & \vdots & \ddots & \vdots & \vdots \\
0 & 0 & \dots & 1 & 0 \\
F_{p_1} & F_{p_2} & \dots & F_{p_{n-1}} & F_{p_n}
\end{pmatrix}.
\end{equation}
Определитель этой матрицы равен:
$$\det D_p G(p^0, x^0) = F_{p_n}(p^0, x^0)$$

Поскольку $F_{p_n} \neq 0$, то по теореме о неявной функции \eqref{th::implicit}
существует единственная гладкая зависимость 
начального градиента $q(y)$ от точек границы $\Gamma$, что доказывает лемму.
\end{proof}

\textbf{4. Следствие о локальной обратимости}
\begin{corollary}
Пусть выполнено условие нехарактеристичности $F_{p_n}(p^0, z^0, x^0) \neq 0$. 
Тогда существуют открытый интервал $I \subset \mathbb{R}$, содержащий ноль 
($0 \in I$), окрестность $W \subset \Gamma \subset \mathbb{R}^{n-1}$ точки 
$x^0$ на границе и открытая окрестность $\Omega \subset \mathbb{R}^n$ точки
$x^0$, такие что для любого $x \in \Omega$ существует единственная пара
переменных $s \in I$ и $y \in W$, для которых выполняется равенство:
\begin{equation}
    x = x(y, s),
\end{equation}
где вектор-функция $x(y, s) \in C^2$ является решением характеристической системы.
\end{corollary}

\begin{proof}
Рассмотрим отображение характеристик $x(y, s)$ с заданным начальным условием
$x(y, 0) = (y, 0)$. Вычислим матрицу Якоби этого преобразования в начальной
точке $(x^0, 0)$.

Частные производные по пространственным координатам вдоль выпрямленной
границы равны:
\begin{equation}
    \frac{\partial x_j}{\partial y_i} = \delta_{ij} \quad \text{при } i = 1, \dots, n-1.
\end{equation}
Производная по параметру вдоль траектории $s$ задается непосредственно самой
характеристической системой \eqref{eq::char_system}:
\begin{equation}
    \frac{\partial x_j}{\partial s} = F_{p_j}(p^0, z^0, x^0).
\end{equation}

Таким образом, определитель матрицы Якоби $D x(x^0, 0)$ совпадает с частной
производной функции $F$ по нормальной компоненте импульса:
\begin{equation}
    \det D x(x^0, 0) = F_{p_n}(p^0, z^0, x^0).
\end{equation}

Так как по условию нехарактеристичности $F_{p_n} \neq 0$, определитель матрицы
Якоби отличен от нуля ($\det D x \neq 0$). Следовательно, по теореме об
обратной функции \eqref{th::implicit}отображение $(y, s) \mapsto x(y, s)$ является локальным диффеоморфизмом класса $C^2$, отображающим область $W \times I$ на некоторую открытую окрестность $\Omega \subset \mathbb{R}^n$ точки $x^0$. 

Это гарантирует, что для каждой точки $x \in \Omega$ существует единственное значение параметра вдоль характеристики $s \in I$ и единственная точка старта на границе $y \in W$, однозначно определяющие проходящую через нее интегральную кривую $x(y,s)$.

\end{proof}

\subsection{Локальная теорема существования}

\begin{theorem}[Локальная теорема существования]
    Пусть $\Gamma$ --- гладкая поверхность, а функции $F$ и $g$
    достаточно гладкие. Если в точке $x^0 \in \Gamma$ выполнено
    условие нехарактеристичности $F_{p_n}(p^0, z^0, x^0) \neq 0$,
    то в некоторой окрестности точки $x^0$ существует единственное
    решение $u(x)$ задачи Коши $F(Du, u, x) = 0$, $u\big|_{\Gamma} = g(x)$.
\end{theorem}

\begin{proof}
    Не умаляя общности, будем считать, что граница локально выпрямлена и
    совпадает с гиперплоскостью $y_n = 0$. Пусть 
    $y = (y_1, \dots, y_{n-1}, 0)$ --- локальные координаты на $\Gamma$ в окрестности $x^0$.
    
    Рассмотрим характеристическую систему, выпущенную из точек границы.
    Решения этой системы зависят от начальной точки $y \in \Gamma$ и параметра вдоль характеристики $s$:
    $x = x(y, s)$, $z = z(y, s)$, $p = p(y, s)$.
    При $s = 0$ выполняются начальные условия:
    \begin{equation}
        x(y, 0) = (y, 0), \quad z(y, 0) = g(y), \quad p(y, 0) = p^0(y).
    \end{equation}
    
    Вычислим матрицу Якоби преобразования $x = x(y, s)$ в точке $s = 0$.
    Дифференцируя по касательным переменным $y_i$, получаем
    $\frac{\partial x_j}{\partial y_i} = \delta_{ij}$ для
    $i, j = 1, \dots, n-1$ и $\frac{\partial x_n}{\partial y_i} = 0$.
    Дифференцируя по параметру $s$, из характеристической системы имеем
    $\frac{\partial x_j}{\partial s} = F_{p_j}(p^0, z^0, x^0)$.
    Тогда определитель матрицы Якоби равен:
    \begin{equation}
        \det D_{y, s} x(y, 0) = F_{p_n}(p^0, z^0, x^0).
    \end{equation}
    
    В силу условия нехарактеристичности $F_{p_n} \neq 0$, определитель отличен от нуля.
    По теореме о локальной обратимости, отображение $x = x(y, s)$
    обратимо в некоторой окрестности точки $x^0$, то есть существуют гладкие функции $y = y(x)$ и $s = s(x)$.
    
    Построим функцию-кандидата на решение:
    \begin{equation}
        u(x) = z(y(x), s(x)).
    \end{equation}
    
    Остается показать, что построенная функция удовлетворяет уравнению
    $F(Du, u, x) = 0$. Введем вспомогательную функцию вдоль характеристик:
    \begin{equation}
        f(y, s) = F(p(y, s), z(y, s), x(y, s)).
    \end{equation}
    При $s = 0$ по построению вектора $p^0(y)$ имеем $f(y, 0) = 0$.
    Дифференцируя $f(y, s)$ по параметру $s$ и используя характеристическую систему, получаем:
    \begin{equation}
        \frac{\partial f(y, s)}{\partial s} = \sum_{j=1}^n \frac{\partial F}{\partial p_j} \dot{p}_j + \frac{\partial F}{\partial z} \dot{z} + \sum_{j=1}^n \frac{\partial F}{\partial x_j} \dot{x}_j.
    \end{equation}
    Подставляя производные из системы характеристик, это выражение тождественно 
    обращается в нуль. Таким образом, $f(y, s) \equiv 0$, что означает выполнение 
    уравнения $F(p, z, x) = 0$ вдоль характеристик. Учитывая соотношение 
    $p_j(y, s) = \frac{\partial u(x(y,s))}{\partial x_j}$, заключаем, что функция $u(x)$ является искомым решением.
\end{proof}
